% Run pdfLaTeX from this directory to resolve the shared inputs.
\makeatletter
\def\input@path{{../../../preamble/}}
\makeatother
\input{preamble.tex}
\input{graphs/highlights.tex}

\geometry{letterpaper,margin=0.55in,bottom=0.67in,footskip=24pt}
\pagestyle{plain}
\linespread{1.10}
\setlength{\parskip}{5pt}
\hypersetup{pdftitle={Grade 10 Term 2 Lesson 1 Classroom Activity: Outcomes, Support and PMFs}}
\newif\ifprintsolutions
\printsolutionstrue % Change to \printsolutionsfalse for the student handout.

\setlength{\arrayrulewidth}{0.55pt}
\setlength{\tabcolsep}{7pt}
\setlength{\extrarowheight}{2pt}
\renewcommand{\arraystretch}{1.20}
\newcolumntype{Y}{>{\centering\arraybackslash}X}
\newcolumntype{L}{>{\raggedright\arraybackslash}X}
\newcommand{\stageheading}[1]{%
    \par\addvspace{0.65\baselineskip}%
    {\centering\large\bfseries #1\par}\nobreak
}
\newcommand{\problemtitle}[2]{%
    \par\addvspace{0.4\baselineskip}%
    \textbf{#1. #2}\par\nobreak
}
\newcommand{\task}[2]{\par\noindent\textbf{Task #1.} #2\par}

\begin{document}
\GTenHighlightsfalse

{\centering
    {\Large\bfseries C1 -- Outcomes, Support and PMFs\par}
    Global Economics \textbar{} Grade 10 \textbar{} Term 2 \textbar{} Lesson 1\par
}

\textbf{C1:} Characterises discrete random variables in economic and financial contexts by identifying their support, probability mass function and possible outcomes.

\stageheading{Worked example}
\problemtitle{1}{Notebook sales}
A school enterprise offers one notebook to each of three visitors. Each visitor buys with probability \(0.4\) and decides independently. Record their decisions in order using \(\mathrm{B}\) for buy and \(\mathrm{N}\) for no purchase. Let \(X\) count the buyers.

\task{1}{Explain what \(X\) counts. List the possible values in buyers and group all ordered
    purchase outcomes by the value of \(X\).}

\task{2}{Build the probability mass function (PMF), \(p_X(x)=P(X=x)\). Check that its masses are
    non-negative and sum to \(1\). State the support and interpret \(p_X(1)\) in context.}

\textbf{Worked reasoning.} The outcomes \(\mathrm{BNN}\), \(\mathrm{NBN}\) and \(\mathrm{NNB}\) differ in who buys, but each gives \(X=1\) buyer. The chance of no purchase is \(1-0.4=0.6\). By independence, each of the three outcomes has probability \(0.4\cdot 0.6^2=0.144\). Only one can occur, so we add their probabilities to get \(p_X(1)\).

\begin{tabularx}{\linewidth}{|c|Y|Y|c|}
    \hline
    \textbf{\(x\) (buyers)} & \textbf{Outcomes giving \(x\)} & \textbf{Probability calculation} & \textbf{\(p_X(x)\)} \\ \hline
    \(0\) & \(\mathrm{NNN}\) & \(0.6^3\) & \(0.216\) \\ \hline
    \(1\) & \(\mathrm{BNN},\,\mathrm{NBN},\,\mathrm{NNB}\) & \(3(0.4\cdot 0.6^2)\) & \(0.432\) \\ \hline
    \(2\) & \(\mathrm{BBN},\,\mathrm{BNB},\,\mathrm{NBB}\) & \(3(0.4^2\cdot 0.6)\) & \(0.288\) \\ \hline
    \(3\) & \(\mathrm{BBB}\) & \(0.4^3\) & \(0.064\) \\ \hline
\end{tabularx}\par
All masses are non-negative and \(0.216+0.432+0.288+0.064=1\). The support consists of the values with positive probability, \(\{0,1,2,3\}\) in buyers. Exactly one visitor buys with probability \(0.432\).

\clearpage
\begingroup
\linespread{1.04}\selectfont
\setlength{\parskip}{3pt}

\stageheading{Guided practice}
\problemtitle{2}{Payment methods and a processing fee}
A shop receives one payment by cash, debit or credit with respective probabilities \(0.25\), \(0.45\) and \(0.30\). Let \(F\) be the processing fee in USD. Cash costs \(0\) USD to process and either card costs \(1\) USD.

\noindent
\begin{minipage}[t]{0.42\linewidth}
    \raggedright
    \task{1}{Explain what \(F\) measures. Map cash, debit and credit to their fees, grouping methods
        that give the same value.}

    \begin{tabularx}{\linewidth}{|c|Y|Y|Y|}
        \hline
        \textbf{Method} & \textbf{Cash} & \textbf{Debit} & \textbf{Credit} \\ \hline
        \textbf{Fee (USD)} & \(0\) & & \\[0.55em] \hline
    \end{tabularx}
\end{minipage}%
\hfill
\begin{minipage}[t]{0.55\linewidth}
    \raggedright
    \setlength{\tabcolsep}{5pt}
    \task{2}{Build the PMF by adding the probabilities of methods with the same fee. Check that the
        masses are non-negative and sum to \(1\). State the support and explain what \(p_F(1)\)
        means for the shop.}

    \begin{tabularx}{\linewidth}{|c|Y|c|}
        \hline
        \textbf{Fee \(f\) (USD)} & \textbf{Probability calculation} & \textbf{Mass \(p_F(f)\)} \\ \hline
        \(0\) & \(P(\text{cash})\) & \\[0.55em] \hline
        \(1\) & & \\[0.55em] \hline
    \end{tabularx}
\end{minipage}

\renewcommand{\arraystretch}{1.10}
\stageheading{Independent practice}
Complete your answers on a separate sheet.

\problemtitle{3}{Delivery bonus}
A delivery arrives early, on time or late with respective probabilities \(0.20\), \(0.50\) and \(0.30\). Exactly one of these outcomes occurs. The delivery bonus \(B\) is \(2\) USD for an early or on-time delivery and \(-3\) USD for a late delivery.

\task{1}{Explain what \(B\) measures and group the three delivery outcomes by their bonus values in
    USD.}

\task{2}{Build the PMF of \(B\) and show how you obtain each mass. Check that the masses are
    non-negative and sum to \(1\). State the support and interpret \(P(B<0)\) in context.}

\problemtitle{4}{Daily insurance claims}
Let \(C\) be the number of insurance claims from a small group of policyholders in one day. The model allows only the counts shown below. The mass \(q\) is unknown.

\begin{tabularx}{\linewidth}{|c|Y|Y|Y|}
    \hline
    \textbf{\(c\) (claims)} & \(0\) & \(1\) & \(2\) \\ \hline
    \textbf{Mass \(p_C(c)\)} & \(0.70\) & \(q\) & \(0.10\) \\ \hline
\end{tabularx}

\task{1}{Explain what \(C\) counts. List the counts in the table and identify the count with an
    unknown mass.}

\task{2}{Find \(q\) to complete the PMF. Check that the masses are non-negative and sum to \(1\).
    State the support, then calculate and interpret the probability of at least one claim in a day.}

\problemtitle{5}{Investment gain}
Let \(G\) be the net gain in USD from a hypothetical one-month investment. Its only possible gain values and their masses are shown below. A student proposes \(q=0.40\).

\begin{tabularx}{\linewidth}{|c|Y|Y|Y|}
    \hline
    \textbf{Net gain \(g\) (USD)} & \(-20\) & \(0\) & \(30\) \\ \hline
    \textbf{Mass \(p_G(g)\)} & \(0.20\) & \(0.50\) & \(q\) \\ \hline
\end{tabularx}

\task{1}{Explain what \(G\) measures. List the gains in USD and identify which listed values
    represent a non-negative gain.}

\task{2}{Test the proposed mass and find the correct value of \(q\). Check that the completed PMF
    has non-negative masses that sum to \(1\). State the support, then calculate and interpret
    \(P(G\geq 0)\).}

\endgroup

\ifprintsolutions
\clearpage
\GTenHighlightstrue
\section*{Solutions}

% Use the G10-T2-C1 legend from graphs/specific_highlights.tex with full cell borders.
\begin{tabularx}{\linewidth}{|l|L|}
    \hline
    \textbf{Highlighted item} & \textbf{What it means} \\ \hline
    \GTenHLOutcome{Outcome} & A possible result of the experiment. \\ \hline
    \GTenHLValue{First variable} & The first variable, its numerical values, and any units. \\ \hline
    \GTenHLSecond{Second variable} & The second variable, its numerical values, and any units. \\ \hline
    \GTenHLThird{Third variable} & The third variable, its numerical values, and any units. \\ \hline
    \GTenHLProbability{Probability} & The chance of an outcome or event, expressed without units. \\ \hline
    \GTenHLParameter{Unknown parameter} & A quantity to determine from the information given. \\ \hline
\end{tabularx}

\problemtitle{2}{Payment methods and a processing fee}
\task{1}{\(\GTenHLValue{F}\) measures the fee for one payment. \GTenHLOutcome{Cash} gives
    \GTenHLValue{0 USD}. \GTenHLOutcome{Debit} and \GTenHLOutcome{credit} both give \GTenHLValue{1
    USD}.}

\task{2}{Cash is the only outcome with a fee of \(0\) USD, so
    \(p_F(\GTenHLValue{0})=\GTenHLProbability{0.25}\). The card methods are mutually exclusive, so
    \(p_F(\GTenHLValue{1})=\GTenHLProbability{0.45}+\GTenHLProbability{0.30}=\GTenHLProbability{0.75}\).
    Both masses are positive and \(\GTenHLProbability{0.25}+\GTenHLProbability{0.75}=1\). The
    support is \(\GTenHLValue{\{0,1\}}\) in USD. The shop pays a \(1\) USD fee with probability
    \(\GTenHLProbability{0.75}\).}

\problemtitle{3}{Delivery bonus}
\task{1}{\(\GTenHLValue{B}\) measures the bonus for one delivery. \GTenHLOutcome{Early} and
    \GTenHLOutcome{on time} both give \GTenHLValue{2 USD}. \GTenHLOutcome{Late} gives
    \GTenHLValue{\(-3\) USD}.}

\task{2}{Add the early and on-time masses to get
    \(p_B(\GTenHLValue{2})=\GTenHLProbability{0.20}+\GTenHLProbability{0.50}=\GTenHLProbability{0.70}\).
    Only a late delivery gives \(-3\) USD, so \(p_B(\GTenHLValue{-3})=\GTenHLProbability{0.30}\).
    Both masses are positive and \(\GTenHLProbability{0.70}+\GTenHLProbability{0.30}=1\). The
    support is \(\GTenHLValue{\{-3,2\}}\) in USD. The bonus is negative with probability
    \(P(B<0)=\GTenHLProbability{0.30}\).}

\problemtitle{4}{Daily insurance claims}
\task{1}{\(\GTenHLValue{C}\) counts insurance claims in one day. The model lists \GTenHLValue{0, 1
    and 2 claims}, with unknown mass \(\GTenHLParameter{q}\) at \(1\) claim.}

\task{2}{The masses must sum to \(1\), so
    \(\GTenHLParameter{q}=1-\GTenHLProbability{0.70}-\GTenHLProbability{0.10}=\GTenHLParameter{0.20}\).
    The completed masses \(\GTenHLProbability{0.70}\), \(\GTenHLProbability{0.20}\) and
    \(\GTenHLProbability{0.10}\) are positive and sum to \(1\). The support is
    \(\GTenHLValue{\{0,1,2\}}\) in claims. At least one claim means \(1\) or \(2\) claims, so
    \(P(C\geq 1)=\GTenHLProbability{0.20}+\GTenHLProbability{0.10}=\GTenHLProbability{0.30}\). There
    is a \(30\%\) chance of a day with at least one claim.}

\problemtitle{5}{Investment gain}
\task{1}{\(\GTenHLValue{G}\) measures the net gain over one month. The listed gains are
    \GTenHLValue{\(-20\), \(0\) and \(30\) USD}. The non-negative gains are \GTenHLValue{0 USD} and
    \GTenHLValue{30 USD}.}

\task{2}{The proposed mass gives
    \(\GTenHLProbability{0.20}+\GTenHLProbability{0.50}+\GTenHLParameter{0.40}=1.10\), so the total
    exceeds \(1\). The correct mass is
    \(\GTenHLParameter{q}=1-\GTenHLProbability{0.20}-\GTenHLProbability{0.50}=\GTenHLParameter{0.30}\).
    The completed masses are positive and
    \(\GTenHLProbability{0.20}+\GTenHLProbability{0.50}+\GTenHLProbability{0.30}=1\). The support is
    \(\GTenHLValue{\{-20,0,30\}}\) in USD. Including zero gives \(P(G\geq
    0)=\GTenHLProbability{0.50}+\GTenHLProbability{0.30}=\GTenHLProbability{0.80}\). The investment
    breaks even or makes a profit with probability \(0.80\).}
\fi

\end{document}
